        Im ersten Schritt sollte Euer Programm etwa so aussehen:

\begin{code}
int main() {
    random_init();

    konsole_loeschen();
    konsole_start();

    int zeile = random_number(1, HOEHE);
    int spalte = random_number(1, BREITE);

    for (int i = 1; i < zeile; i++) {
        konsole_zeichen('\n');
    }
    for (int s = 1; s < spalte; s++) {
        konsole_zeichen(' ');
    }
    konsole_zeichen('*');

    return 0;
}
\end{code}
    Der Ablauf im Einzelnen ist:
    \begin{enumerate}
    \item Initialisierung des Zufallszahlengenerators.
    \item Konsole löschen und in die erste Zeile und Spalte gehen.
    \item Berechnen der Zufallszahlen für Zeile und Spalte, in der der Stern ausgegeben werden soll.
    \item Schleife, um in die richtige Zeile zu gehen.
    \item Schleife, um in die richtige Spalte zu gehen.
    \item Ausgabe des Sterns.
    \end{enumerate}

    \noindent
    Im zweiten Schritt folgt dann:
\begin{code}
int main() {
    random_init();

    konsole_loeschen();
    while (1) {
        konsole_start();

        int zeile = random_number(1, HOEHE);
        int spalte = random_number(1, BREITE);

        for (int i = 1; i < zeile; i++) {
            konsole_zeichen('\n');
        }
        for (int s = 1; s < spalte; s++) {
            konsole_zeichen(' ');
        }
        konsole_zeichen('*');

        usleep(100000);
    }

    return 0;
}
\end{code}
    Hier kommt nur das \codeinline{while} und \codeinline{usleep} hinzu.
